\chapter{部分積分}\label{章:部分積分}

前章では,合成関数の微分を逆向きに読んで,複雑なまとまりを一つの変数に換えた.たとえば $x(x^2+1)^{20}$ には,\,$x^2+1$ の微分に当たる $2x$ が,定数倍を除いて含まれていた.ところが $x\e^x$ では,\,$x$ と $\e^x$ のどちらかを新しい変数にしても,もう片方と $\dd x$ が一緒に簡単な形へ変わるわけではない.

ここで,積の微分を思い出す.\,$x\e^x$ を微分すると,\,$x\e^x$ のほかに $\e^x$ が現れる.したがって,\,$x\e^x$ 自体は求める原始関数ではないが,余分に現れる $\e^x$ の原始関数を引けばよさそうである.このように,積を微分したときの二つの項を使って,ある積分を別の積分へ換えるのが部分積分である.

\section{部分積分の公式}\label{節:部分積分の導出}

一般的な形は
\[
 \int u(x)v'(x)\dd x
\]
で与えられる.ここで $u$ は微分すると簡単になる因子,\,$v'$ は原始関数 $v$ が分かる因子である.この積分を,\,$u$ を微分し $v'$ を積分した積 $u'v$ の積分に帰着させることを目標とする.

2つの関数 $u(x),v(x)$ が微分できるとき,積の微分は
\[
 \{u(x)v(x)\}'=u'(x)v(x)+u(x)v'(x)
\]
である.求めたいのは右辺の第2項の積分なので,第1項を移項して
\[
 u(x)v'(x)=\{u(x)v(x)\}'-u'(x)v(x)
\]
と書き直してやる.両辺の原始関数を考えると,右辺の第1項からは積 $uv$ が,第2項からはもう1つの積分が現れる.

\begin{theorem}{部分積分の公式}
$u,v$ がある区間で連続微分可能ならば,その区間で
\[
 \int u(x)v'(x)\dd x
 =u(x)v(x)-\int u'(x)v(x)\dd x
\]
が成り立つ.不定積分は任意定数を含めた原始関数の族として読む.特に,同じ区間の $a,b$ に対して
\[
 \int_a^b u(x)v'(x)\dd x
 =[u(x)v(x)]_a^b-\int_a^b u'(x)v(x)\dd x
\]
が成り立つ.
\end{theorem}

定積分の公式は,積の微分を $a$ から $b$ まで積分し,基本定理を使ったものである.\,$[uv]_a^b$ は $u(b)v(b)-u(a)v(a)$ を表す.積の微分のうち一方の項だけを移項したので,この端点の項を省くことはできない.

公式の左辺では $u$ をそのまま,\,$v'$ を微分された形で書いている.右辺へ移ると,\,$u$ は $u'$ へ微分され,\,$v'$ は原始関数 $v$ に変わる.したがって実際に考えることは,「どの因子を微分し,どの因子を積分すると,残る積 $u'v$ が扱いやすくなるか」である.

ときには $\dd v=v'(x)\dd x$,\,$\dd u=u'(x)\dd x$ とまとめて,
\[
 \int u\dd v=uv-\int v\dd u
\]
とも書く.これは同じ公式の短い表記であり,\,$u$ と $v$ を関係のない変数として自由に動かしているのではない.本章では,微分する因子と積分する因子を確かめやすいように,主に $u,v'$ を書いて計算する.

ここで1つ断りをいれておこう.本章では,不定積分の $C$ は任意定数を表し,現れる関数は考える区間で連続微分可能とする.対数を含む式は真数が正になる範囲で考える.以後この条件を本文では繰り返さないが,答案等に書く際は常に確かめること.

\begin{bookpoint}
\textbf{POINT}\par
微分すると簡単になる因子を $u$,原始関数が分かる因子を $v'$ と選び,\,$\int uv'\dd x=uv-\int u'v\dd x$ で残りの積分 $\int u'v\dd x$ へ移す.
\end{bookpoint}

\begin{example}{一次式と指数関数}
$\int x\e^x\dd x$ を求めよ.
\end{example}

\paragraphbf{方針}
$x$ を微分して $1$ にし,\,$\e^x$ を積分する.残る積分は $\int\e^x\dd x$ である.

\paragraphbf{解答}
$(x-1)\e^x+C$

\paragraphbf{解法}
$x$ は微分すると $1$ になり,\,$\e^x$ は積分しても $\e^x$ のままである.そこで
\[
 u=x,\qquad v'=\e^x,
 \qquad u'=1,\qquad v=\e^x
\]
と選ぶ.部分積分の公式より,
\[
\begin{aligned}
 \int x\e^x\dd x
 &=x\e^x-\int 1\cdot\e^x\dd x\\
 &=x\e^x-\e^x+C\\
 &=(x-1)\e^x+C
\end{aligned}
\]
を得る.答えを微分すると
\[
 \frac{\dd}{\dd x}\{(x-1)\e^x\}
 =\e^x+(x-1)\e^x
 =x\e^x
\]
となり,確かにもとの関数に戻る.

\paragraphbf{引いた項の意味}
章の初めで述べたように,積 $x\e^x$ を作って終わりにすると,その微分は $x\e^x+\e^x$ となる.例題の答えの $-\e^x$ は,この余分な $\e^x$ を取り除くために引いた項である.部分積分の公式の $-\int u'v\dd x$ は,一般にこの「余分な項」の原始関数を表しているのだ.

\paragraphbf{因子を逆に選んだ場合}
逆に $u=\e^x$,\,$v'=x$ と選ぶことも,公式上は正しい.その場合には
\[
 \int x\e^x\dd x
 =\frac{x^2}{2}\e^x-\frac12\int x^2\e^x\dd x
\]
となる.ところが,もとの一次式が二次式になり,残った積分はむしろ手間が増えている.公式を使えることと,その選択で計算が進むことは別である.右辺に何が残るかを一度書き下してから,因子の選び方を判断すればよい.
%
\section{因子の選び方}\label{節:部分積分の因子選択}

前節の例題では,一次式 $x$ を微分し,指数関数 $\e^x$ を積分した.逆に選ぶと計算が重くなることも見た.では一般に,どちらの因子を微分し,どちらを積分すればよいのだろうか.目標は,残る積分 $\int u'v\dd x$ が既知の積分になるように選ぶことである.

まず,積分する側には条件がある.原始関数 $v$ を,その時点で求められなければならない.指数関数や三角関数は,積分しても同じ種類の関数のままなので,積分する側に選びやすい.一方,整式は微分すると次数が下がるため,微分する側の候補になる.

ただし,この順序は絶対の規則ではない.対数 $\log x$ の原始関数はまだ求めていないので,\,$\log x$ は積分する側には置けない.そこで $\log x$ を微分する側に選んでやる.\,$n\neq-1$ とし,
\[
 \int x^n\log x\dd x
\]
を考えよう.\,$u=\log x$,\,$v'=x^n$ と選ぶと,
\[
 u'=\frac1x,\qquad v=\frac{x^{n+1}}{n+1}
\]
である.部分積分の公式より,
\[
\begin{aligned}
 \int x^n\log x\dd x
 &=\frac{x^{n+1}}{n+1}\log x-\int\frac{x^{n+1}}{n+1}\cdot\frac1x\dd x\\
 &=\frac{x^{n+1}}{n+1}\log x-\frac{1}{n+1}\int x^n\dd x\\
 &=\frac{x^{n+1}}{n+1}\log x-\frac{x^{n+1}}{(n+1)^2}+C
\end{aligned}
\]
を得る.対数を微分して得た $\frac1x$ が,整式を積分して得た $x^{n+1}$ と約分され,残る積分が整式になったのだ.ここでは整式 $x^n$ の方を積分する側に選んでいる.関数の名前だけで順序を決めず,\,$u'v$ を具体的に見ることが肝である.

特に $n=0$ のときは $x^0=1$ なので,被積分関数は $\log x$ だけである.見た目が積でなくても,\,$\log x\cdot1$ と $1$ を掛けた積と考えれば,積分する側に既知の関数 $1$ を置けるのである.

\begin{bookpoint}
\textbf{POINT}\par
整式は微分する側,指数関数・三角関数は積分する側に選ぶ.対数は微分する側に選び,\,$\frac1x$ を残りの因子と約分する.
\end{bookpoint}

\begin{example}{1つの関数を積と見る}
$x>0$ において,\,$\int\log x\dd x$ を求めよ.
\end{example}

\paragraphbf{方針}
$\log x=\log x\cdot1$ と見て,\,$\log x$ を微分し,\,$1$ を積分する.

\paragraphbf{解答}
$x\log x-x+C$

\paragraphbf{解法}
$\log x$ の微分は $\frac1x$ であり,これに $1$ の原始関数 $x$ が掛かれば積は $1$ になる.そこで
\[
 u=\log x,\qquad v'=1,
 \qquad u'=\frac1x,\qquad v=x
\]
と選ぶと,
\[
\begin{aligned}
 \int\log x\dd x
 &=x\log x-\int x\cdot\frac1x\dd x\\
 &=x\log x-\int 1\dd x\\
 &=x\log x-x+C
\end{aligned}
\]
となる.答えを微分すると,
\[
 \frac{\dd}{\dd x}(x\log x-x)
 =\log x+x\cdot\frac1x-1
 =\log x
\]
である.

\paragraphbf{置換を先に考える場合}
上の一般形で除いた $n=-1$ の場合,すなわち $\int\frac{\log x}{x}\dd x$ では,\,$v=\frac{x^{n+1}}{n+1}$ が作れない.しかしこの場合は,前章で見たとおり $t=\log x$ と置くほうが短い.\,$\frac1x\dd x=\dd t$ となり,ただちに $\int t\dd t$ へ変わるからである.対数を見たら部分積分と決めるのではなく,対数の微分 $\frac1x$ がすでに残りの部分にあるかを先に調べよう.

積分する側の条件についても同じことがいえる.たとえば $\int x\e^{x^2}\dd x$ で,\,$x$ を微分しようとして $v'=\e^{x^2}$ と選ぶと,そこで原始関数が分からなくなる.しかし $t=x^2$ と置けば,\,$x\dd x=\frac12\dd t$ なので,
\[
 \int x\e^{x^2}\dd x
 =\frac12\int\e^t\dd t
 =\frac12\e^{x^2}+C
\]
である.部分積分を試す前に,置換に必要な微分の因子が含まれていないかを見ると,遠回りを避けられる.
%
\section{部分積分の反復}\label{節:部分積分の反復}

前節までは,一度の部分積分で既知の積分に到達した.では,整式の次数が高い場合はどうなるだろうか.一般的な形は,\,$P(x)$ を $n$ 次の整式,\,$a\neq0$ として\footnote{$a=0$ なら $\int P(x)\dd x$ であり,部分積分は要らない.}
\[
 \int P(x)\e^{ax}\dd x,\qquad
 \int P(x)\cos ax\dd x,\qquad
 \int P(x)\sin ax\dd x
\]
で与えられる.整式を微分する側に選ぶと,一度の部分積分で整式の次数が1つ下がる.これを繰り返して整式を定数にし,\,$\int\e^{ax}\dd x$ などの既知の積分に帰着させることを目標とする.

指数関数の場合を書き下してみよう.\,$u=P(x)$,\,$v'=\e^{ax}$ と選ぶと,\,$v=\frac1a\e^{ax}$ だから
\[
 \int P(x)\e^{ax}\dd x
 =\frac1aP(x)\e^{ax}-\frac1a\int P'(x)\e^{ax}\dd x
\]
となる.右辺の積分は,整式の次数が1つ下がっただけで,もとと同じ形である.そこで同じ部分積分をもう一度施してやれば,
\[
 \int P(x)\e^{ax}\dd x
 =\frac1aP(x)\e^{ax}-\frac1{a^2}P'(x)\e^{ax}
  +\frac1{a^2}\int P''(x)\e^{ax}\dd x
\]
となる.ここで,\,$P$ を $k$ 回微分した関数を\textbf{第$k$次導関数}といい,\,$P^{(k)}$ と書く.\,$P^{(1)}=P'$,\,$P^{(2)}=P''$ である.\,$n$ 次の整式は微分するたびに次数が1つ下がるので,\,$P,P',P'',\ldots$ と微分を繰り返すと,第$n$次導関数 $P^{(n)}$ は定数になる.したがって $n$ 回目の部分積分の後に残る積分は $\int\e^{ax}\dd x$ の定数倍であり,ここで計算が終わるのだ.まとめると
\[
 \int P(x)\e^{ax}\dd x
 =\e^{ax}\left\{\frac{P(x)}{a}-\frac{P'(x)}{a^2}+\frac{P''(x)}{a^3}-\cdots
  +(-1)^n\frac{P^{(n)}(x)}{a^{n+1}}\right\}+C
\]
を得る.

三角関数の場合も同じ手順で進む.異なるのは,\,$\cos ax$ と $\sin ax$ を積分するたびに $\frac1a$ が掛かるだけでなく,正弦と余弦が入れ替わり,符号も変わることである.次数が下がることと,係数・符号が変わることを別々に確認すればよい.

\begin{bookpoint}
\textbf{POINT}\par
整式を微分する側,指数関数・三角関数を積分する側に選び,整式が定数になるまで部分積分を繰り返す.
\end{bookpoint}

\begin{example}{部分積分の反復}
$\int x^2\cos x\dd x$ を求めよ.
\end{example}

\paragraphbf{方針}
$x^2$ を微分し,\,$\cos x$ を積分する.残る $\int x\sin x\dd x$ でもう一度,整式を微分する部分積分を行う.

\paragraphbf{解答}
$x^2\sin x+2x\cos x-2\sin x+C$

\paragraphbf{解法}
$u=x^2$,\,$v'=\cos x$ と選ぶ.\,$u'=2x$,\,$v=\sin x$ なので,
\[
 \int x^2\cos x\dd x
 =x^2\sin x-2\int x\sin x\dd x
\]
となる.残った積分では,\,$x$ を微分し,\,$\sin x$ を積分する.\,$\sin x$ の原始関数は $-\cos x$ だから,
\[
\begin{aligned}
 \int x\sin x\dd x
 &=x(-\cos x)-\int1\cdot(-\cos x)\dd x\\
 &=-x\cos x+\int\cos x\dd x\\
 &=-x\cos x+\sin x+C
\end{aligned}
\]
である.これをもとの式へ戻して,任意定数をまとめると
\[
\begin{aligned}
 \int x^2\cos x\dd x
 &=x^2\sin x-2(-x\cos x+\sin x)+C\\
 &=x^2\sin x+2x\cos x-2\sin x+C
\end{aligned}
\]
を得る.2回目の積分で現れた負号だけでなく,その式全体に最初の $-2$ が掛かることに注意する.途中の積分を独立に計算したときは,戻す先の係数を残しておくと,符号の取り違えを防げる.

\paragraphbf{答えの微分による確認}
答えを微分すると,途中で現れた項が対になって消える.
\[
\begin{aligned}
 &\frac{\dd}{\dd x}
  \{x^2\sin x+2x\cos x-2\sin x\}\\
 &\qquad=2x\sin x+x^2\cos x
       +2\cos x-2x\sin x-2\cos x\\
 &\qquad=x^2\cos x
\end{aligned}
\]
$x^2\sin x$ の微分で生まれた余分な項 $2x\sin x$ を,\,$2x\cos x$ の微分の $-2x\sin x$ が打ち消す.その際に生まれた $2\cos x$ を,さらに $-2\sin x$ の微分が打ち消す.部分積分で補った項は,このように積の微分で生まれた余分な項を順に打ち消しているのである.

\paragraphbf{積分の漸化式との関係}
次数を文字のまま残しておくこともできる.\,$I_n=\int x^n\e^{ax}\dd x$ と置き,\,$x^n$ を微分する部分積分を1回だけ行えば,
\[
 I_n=\frac1ax^n\e^{ax}-\frac naI_{n-1}
\]
となる.これは $I_n$ を $I_{n-1}$ で表す漸化式である.このように,部分積分で次数を1つ下げる操作は,積分を項とする数列の漸化式として書き下せる.\chapterref{章:積分と数列}では,この形の式を使って定積分の値を求める.
%
\section{循環する積分}\label{節:部分積分の循環}

指数関数と三角関数の積では,整式のように微分を繰り返して $0$ になる因子がない.では,部分積分を繰り返すとどうなるだろうか.一般的な形は,\,$a\neq0$ として\footnote{$a=0$ なら $\int\sin bx\dd x$ であり,部分積分は要らない.}
\[
 \int\e^{ax}\sin bx\dd x,\qquad
 \int\e^{ax}\cos bx\dd x
\]
で与えられる.正弦と余弦は微分によって互いに移り変わるので,2回の部分積分で,もとの積分が再び現れる.そこで計算を失敗と決めず,もとの積分を未知数として含む等式を作り,それを方程式として解くことを目標とする.

積分定数に惑わされずに等式を整理するため,まず
\[
 I(x)=\int_0^x\e^{at}\sin bt\dd t,\qquad
 J(x)=\int_0^x\e^{at}\cos bt\dd t
\]
と定める.\,$I,J$ は $x$ の関数であり,基本定理から $I'(x)=\e^{ax}\sin bx$ である.したがって $I(x)$ が求まれば,それに任意定数を加えて不定積分の答えにできる.

$I(x)$ で $\sin bt$ を微分し,\,$\e^{at}$ を積分すると,
\[
\begin{aligned}
 I(x)
 &=\left[\frac1a\e^{at}\sin bt\right]_0^x-\frac ba\int_0^x\e^{at}\cos bt\dd t\\
 &=\frac1a\e^{ax}\sin bx-\frac baJ(x)
\end{aligned}
\]
となる.続いて $J(x)$ でも $\e^{at}$ を積分する側に選ぶ.\,$\cos bt$ の微分は $-b\sin bt$ だから,
\[
\begin{aligned}
 J(x)
 &=\left[\frac1a\e^{at}\cos bt\right]_0^x+\frac ba\int_0^x\e^{at}\sin bt\dd t\\
 &=\frac1a\left(\e^{ax}\cos bx-1\right)+\frac baI(x)
\end{aligned}
\]
となる.ここで $I(x)$ が戻ってきた.第2式を第1式へ代入すると
\[
 I(x)
 =\frac1a\e^{ax}\sin bx-\frac b{a^2}\left(\e^{ax}\cos bx-1\right)-\frac{b^2}{a^2}I(x)
\]
となる.右辺の $I(x)$ を左辺へ移項し,両辺に $a^2$ を掛ければ
\[
 (a^2+b^2)I(x)=\e^{ax}(a\sin bx-b\cos bx)+b
\]
を得る.\,$a^2+b^2>0$ なので両辺をこれで割ることができ,
\[
 I(x)=\frac{\e^{ax}(a\sin bx-b\cos bx)}{a^2+b^2}+\frac{b}{a^2+b^2}
\]
となる.最後の定数を任意定数にまとめれば,不定積分が求まったのだ.

\begin{bookpoint}
\textbf{POINT}\par
指数関数を2回とも積分する側に選び,もとの積分 $I$ が戻った等式を $I$ についての方程式として解く.
\end{bookpoint}

\begin{example}{指数関数と正弦の積}
$\int\e^x\sin x\dd x$ を求めよ.
\end{example}

\paragraphbf{方針}
$I(x)=\int_0^x\e^t\sin t\dd t$,\,$J(x)=\int_0^x\e^t\cos t\dd t$ と置き,\,$\e^t$ を積分する部分積分を2回行って $I(x)$ についての方程式を作る.

\paragraphbf{解答}
$\frac12\e^x(\sin x-\cos x)+C$

\paragraphbf{解法}
$I(x)$ で $\sin t$ を微分し,\,$\e^t$ を積分すると,
\[
\begin{aligned}
 I(x)
 &=[\e^t\sin t]_0^x-\int_0^x\e^t\cos t\dd t\\
 &=\e^x\sin x-J(x)
\end{aligned}
\]
となる.\,$J(x)$ でも $\e^t$ を積分する側に選ぶと,
\[
\begin{aligned}
 J(x)
 &=[\e^t\cos t]_0^x-\int_0^x\e^t(-\sin t)\dd t\\
 &=\e^x\cos x-1+I(x)
\end{aligned}
\]
である.第2式を第1式へ代入すると
\[
 I(x)
 =\e^x\sin x-\{\e^x\cos x-1+I(x)\}
\]
となる.\,$I(x)$ を移項して
\[
 2I(x)=\e^x(\sin x-\cos x)+1
\]
を得る.よって
\[
 I(x)=\frac12\e^x(\sin x-\cos x)+\frac12
\]
である.最後の $\frac12$ も任意定数へまとめれば,
\[
 \int\e^x\sin x\dd x
 =\frac12\e^x(\sin x-\cos x)+C
\]
となる.これは上の一般形で $a=b=1$ とした場合に一致する.微分による確認は
\[
\begin{aligned}
 &\frac{\dd}{\dd x}\left\{\frac12\e^x(\sin x-\cos x)\right\}\\
 &\qquad=\frac12\e^x(\sin x-\cos x)
 +\frac12\e^x(\cos x+\sin x)\\
 &\qquad=\e^x\sin x
\end{aligned}
\]
である.

\paragraphbf{同じ積分が戻っても解けない場合}
ここで役立ったのは,「同じ積分が戻る」という事実だけではない.右辺へ $-I(x)$ として戻ったので,移項すると $2I(x)$ が得られたことが肝である.一般形でも,戻ってきた $I(x)$ の係数は $-\frac{b^2}{a^2}$ であり,移項した後の係数 $1+\frac{b^2}{a^2}$ は $0$ にならなかった.

一方,\,$J(x)$ の部分積分で,\,$\e^t$ を微分する側に選び直したとしよう.すると1回目の計算を逆向きにたどることになり,\,$I(x)=I(x)$ という等式に戻ってしまう.これは正しい等式だが,\,$I(x)$ についての新しい情報はない.循環した後には,何倍の $I(x)$ が残ったのかを確かめるとよい.

\paragraphbf{不定積分のまま計算する場合}
不定積分のまま同じ計算を進めることもできる.その場合は,途中の等式が任意定数の違いを許していることを忘れず,最後に1つの $C$ にまとめる.定積分として $I,J$ を定めたのは,等式の両辺を同じ関数として固定し,移項の意味をはっきりさせるためである.
%
\section{置換との併用}\label{節:置換と部分積分}

積分法は,問題ごとに1つだけ選ぶものではない.前章でも,根号を変数に換えた後で割り算を使った.では,置換と部分積分を組み合わせるとどうなるだろうか.ここでは
\[
 \int f(\log x)\dd x
\]
の形の積分を考える.対数の中身を1つの変数にする置換をすれば積が現れ,それを部分積分で計算できる.

まず,\,$f$ の中にある $\log x$ を1つの変数 $t$ にしたい.しかし,もとの積分には $\log x$ の微分 $\frac1x$ がないので,\,$f(t)\dd t$ とそのまま置き換えることはできない.そこで $t=\log x$ を $x=\e^t$ と解き直してやる.\,$\log x$ は狭義単調増加なので,逆向きの対応 $x=\e^t$ も1つに決まる.これを微分すると
\[
 \dd x=\e^t\dd t
\]
である.この微分の係数も含めて換えると,
\[
 \int f(\log x)\dd x=\int f(t)\e^t\dd t
\]
となる.置換によって $\dd x$ から $\e^t$ という因子が現れ,右辺は $f(t)$ と $\e^t$ の積になった.\,$\e^t$ は積分しても $\e^t$ のままなので,積分する側にうってつけである.\,$u=f(t)$,\,$v'=\e^t$ として部分積分すれば,
\[
 \int f(t)\e^t\dd t
 =f(t)\e^t-\int f'(t)\e^t\dd t
\]
となる.\,$f$ が整式や三角関数なら,前節までの方法で右辺の積分を求められる.最後に $t=\log x$ を代入して,もとの変数へ戻せばよい.

\begin{bookpoint}
\textbf{POINT}\par
$t=\log x$ と置いて $\dd x=\e^t\dd t$ とし,現れた積 $f(t)\e^t$ を部分積分してから $t=\log x$ に戻す.
\end{bookpoint}

\begin{example}{対数と三角関数}
$x>0$ において,\,$\int\sin(\log x)\dd x$ を求めよ.
\end{example}

\paragraphbf{方針}
$t=\log x$ と置くと $\int\e^t\sin t\dd t$ になる.これは前節の例題の形である.

\paragraphbf{解答}
$\frac x2\{\sin(\log x)-\cos(\log x)\}+C$

\paragraphbf{解法}
$t=\log x$ と置くと,\,$x>0$ だから $x=\e^t$ であり,\,$\dd x=\e^t\dd t$ となる.したがって
\[
 \int\sin(\log x)\dd x
 =\int\e^t\sin t\dd t
\]
である.前節の例題の結果を使うと,
\[
\begin{aligned}
 \int\sin(\log x)\dd x
 &=\frac12\e^t(\sin t-\cos t)+C\\
 &=\frac x2\{\sin(\log x)-\cos(\log x)\}+C
\end{aligned}
\]
を得る.最後の行では $\e^t=x$,\,$t=\log x$ を代入した.

検算では,外にある $x$ と,内側の $\log x$ の両方を微分する.
\[
\begin{aligned}
 &\frac{\dd}{\dd x}
 \left[\frac x2\{\sin(\log x)-\cos(\log x)\}\right]\\
 &\qquad=\frac12\{\sin(\log x)-\cos(\log x)\}
   +\frac x2\cdot\frac1x
     \{\cos(\log x)+\sin(\log x)\}\\
 &\qquad=\sin(\log x)
\end{aligned}
\]

\paragraphbf{手法の役割}
$t=\log x$ と置くだけでは計算は終わらず,置換後の積を部分積分で求め,最後に元の変数へ戻した.置換は合成関数の中身を1つの変数にする役割を,部分積分は $\e^t$ との積を処理する役割を担っている.途中の式の形に合わせて,次に使う方法を選べばよいのである.
%
\section{根号を含む積分}\label{節:平方根と部分積分}

$\sqrt{1+x^2}$ も,そのままでは基本公式の形ではない.前章の三角置換から $x=\tan\theta$ を考えることはできるが,ここでは部分積分を使う別の道を調べよう.一般的な形は,\,$A>0$ として\footnote{$A<0$ のときは $x+\sqrt{x^2+A}$ が負になる範囲があり,対数の中身に絶対値が必要になる.本節では $A>0$ に限る.}
\[
 \int\sqrt{x^2+A}\dd x
\]
で与えられる.循環する積分と同じく,部分積分でもとの積分を含む等式を作り,それを解くことを目標とする.

手がかりは,根号を微分すると $\frac{x}{\sqrt{x^2+A}}$ になり,さらに $x$ を掛ければ分子に $x^2$ ができることである.この $x^2$ を $(x^2+A)-A$ と分けると,もとの根号が再び現れる.

$u=\sqrt{x^2+A}$,\,$v'=1$ と選ぶ.\,$u'=\frac{x}{\sqrt{x^2+A}}$,\,$v=x$ だから,部分積分の公式より
\[
 \int\sqrt{x^2+A}\dd x
 =x\sqrt{x^2+A}-\int\frac{x^2}{\sqrt{x^2+A}}\dd x
\]
を得る.右辺の積分を見ただけでは簡単になったか分かりにくいが,分子を分けると
\[
\begin{aligned}
 \frac{x^2}{\sqrt{x^2+A}}
 &=\frac{(x^2+A)-A}{\sqrt{x^2+A}}\\
 &=\sqrt{x^2+A}-\frac{A}{\sqrt{x^2+A}}
\end{aligned}
\]
となる.したがって,もとの積分を左辺へ移項すれば,任意定数を除いて
\[
 2\int\sqrt{x^2+A}\dd x
 =x\sqrt{x^2+A}+A\int\frac1{\sqrt{x^2+A}}\dd x
\]
という関係が得られる.未知の積分がもう一度現れたが,その係数をまとめることができた.残る課題は,根号の逆数の積分である.

この原始関数は少々発想しにくいので,結果からお見せしよう.
\[
 \int\frac1{\sqrt{x^2+A}}\dd x
 =\log\left(x+\sqrt{x^2+A}\right)+C
\]
である.実際,\,$r(x)=\sqrt{x^2+A}$ と略記すると $r'(x)=\frac{x}{r(x)}$ だから,
\[
\begin{aligned}
 \frac{\dd}{\dd x}\log\{x+r(x)\}
 &=\cfrac{1+\cfrac{x}{r(x)}}{x+r(x)}\\
 &=\frac{1}{x+r(x)}\cdot\frac{r(x)+x}{r(x)}\\
 &=\frac1{r(x)}
\end{aligned}
\]
となる.また $r(x)=\sqrt{x^2+A}>|x|$ なので,\,$x+r(x)>0$ はすべての実数 $x$ で成り立ち,対数の中身は正である.これを先ほどの関係へ戻すと,
\[
 \int\sqrt{x^2+A}\dd x
 =\frac12\left\{x\sqrt{x^2+A}
       +A\log\left(x+\sqrt{x^2+A}\right)\right\}+C
\]
を得る.

\begin{bookpoint}
\textbf{POINT}\par
根号を $u$,\,$1$ を $v'$ として部分積分し,分子を $x^2=(x^2+A)-A$ と分けてもとの積分を移項する.残る根号の逆数は $\log\left(x+\sqrt{x^2+A}\right)$ で積分する.
\end{bookpoint}

\begin{example}{根号の部分積分}
$\int\sqrt{1+x^2}\dd x$ を求めよ.
\end{example}

\paragraphbf{方針}
$\sqrt{1+x^2}$ を微分し,\,$1$ を積分する.分子を $x^2=(1+x^2)-1$ と分け,もとの積分を移項する.

\paragraphbf{解答}
$\frac12\left\{x\sqrt{1+x^2}+\log\left(x+\sqrt{1+x^2}\right)\right\}+C$

\paragraphbf{解法}
$u=\sqrt{1+x^2}$,\,$v'=1$ と選ぶ.\,$u'=\frac{x}{\sqrt{1+x^2}}$,\,$v=x$ だから,
\[
 \int\sqrt{1+x^2}\dd x
 =x\sqrt{1+x^2}-\int\frac{x^2}{\sqrt{1+x^2}}\dd x
\]
である.分子を分けると
\[
 \frac{x^2}{\sqrt{1+x^2}}
 =\sqrt{1+x^2}-\frac1{\sqrt{1+x^2}}
\]
となるので,もとの積分を移項して
\[
 2\int\sqrt{1+x^2}\dd x
 =x\sqrt{1+x^2}+\int\frac1{\sqrt{1+x^2}}\dd x
\]
を得る.右辺の積分は $\log\left(x+\sqrt{1+x^2}\right)+C$ であり,\,$x+\sqrt{1+x^2}>0$ がすべての実数 $x$ で成り立つ.よって
\[
 \int\sqrt{1+x^2}\dd x
 =\frac12\left\{x\sqrt{1+x^2}
       +\log\left(x+\sqrt{1+x^2}\right)\right\}+C
\]
である.最後に答え全体を微分すると,
\[
\begin{aligned}
 &\frac12\left\{\sqrt{1+x^2}
       +\frac{x^2}{\sqrt{1+x^2}}
       +\frac1{\sqrt{1+x^2}}\right\}\\
 &\qquad=\frac12\{\sqrt{1+x^2}+\sqrt{1+x^2}\}\\
 &\qquad=\sqrt{1+x^2}
\end{aligned}
\]
となる.この確認により,途中で不定積分を移項した計算も,最終的な原始関数として正しいことが分かる.

\paragraphbf{対数の式の発想}
$\log\left(x+\sqrt{x^2+A}\right)$ という式は,どこから出てきたのだろうか.\,$\frac1{r(x)}$ を積分したいので,微分すると分母に $r(x)$ が現れる式を探す.\,$r'(x)=\frac{x}{r(x)}$ はその候補だが,分子に $x$ が残る.そこで $x'=1$ も,分母を $r(x)$ にして $\frac{r(x)}{r(x)}$ と並べてみよう.この2つを足せば,
\[
 \{x+r(x)\}'=\frac{x+r(x)}{r(x)}
\]
となり,微分する前の $x+r(x)$ が分子に現れる.ここで対数の微分 $\{\log g(x)\}'=\frac{g'(x)}{g(x)}$ を思い出せば,\,$g(x)=x+r(x)$ として分子と分母の $x+r(x)$ が約分され,\,$\frac1{r(x)}$ だけが残るのである.

\paragraphbf{他の置換との関係}
同じ積分は置換でも求められる.\,$x=\tan\theta$ と置くと $\int\sqrt{1+x^2}\dd x$ は $\int\frac{\dd\theta}{\cos^3\theta}$ となり,ここでも部分積分が必要になる.また,本節で現れた $x+\sqrt{1+x^2}$ そのものを $t$ と置く置換もある.これらは\chapterref{章:置換の背景}で双曲線関数とともに扱い,本節の答えと比較する.
%
\section{定積分の部分積分}\label{節:定積分の部分積分}

定積分を求めるとき,不定積分を求めてから上下端を代入してもよい.しかし最初から定積分の公式
\[
 \int_a^b u(x)v'(x)\dd x=[u(x)v(x)]_a^b-\int_a^b u'(x)v(x)\dd x
\]
を使えば,端点で消える項がすぐに分かる.本節では,定積分を端点の項 $[uv]_a^b$ と,残りの定積分 $\int_a^b u'v\dd x$ の2つに分けて計算する.

端点の項は
\[
 [u(x)v(x)]_a^b=u(b)v(b)-u(a)v(a)
\]
であり,数を代入するだけで値が決まる.特に,端点で $u$ か $v$ のどちらかが $0$ になれば,その端点の寄与はなくなる.したがって,残る積分が簡単になるかに加えて,\,$u(a),u(b),v(a),v(b)$ の値も計算の見通しに関わる.端点の項を先に計算しておけば,あとは残りの定積分に集中できる.

\begin{bookpoint}
\textbf{POINT}\par
端点の項 $[uv]_a^b$ を先に計算し,残る $\int_a^b u'v\dd x$ を求めて引く.
\end{bookpoint}

\begin{example}{端点の値}
$\int_0^{\frac\pi2}x\sin x\dd x$ を求めよ.
\end{example}

\paragraphbf{方針}
$x$ を微分し,\,$\sin x$ を積分する.端点の項は上端でも下端でも $0$ になる.

\paragraphbf{解答}
$1$

\paragraphbf{解法}
$u=x$,\,$v'=\sin x$ と選ぶと,\,$u'=1$,\,$v=-\cos x$ なので,
\[
\begin{aligned}
 \int_0^{\frac\pi2}x\sin x\dd x
 &=[-x\cos x]_0^{\frac\pi2}
   +\int_0^{\frac\pi2}\cos x\dd x\\
 &=0+[\sin x]_0^{\frac\pi2}\\
 &=1
\end{aligned}
\]
となる.上端では $\cos\frac\pi2=0$,下端では $x=0$ のため,端点の項はともに $0$ になった.2つの端点で $0$ になる理由は異なり,一方だけを見て項全体を消しているわけではない.

\paragraphbf{値の符号による確認}
積分される関数 $x\sin x$ はこの区間で非負だから,求めた値が正であることは自然である.仮に $-1$ になったら,\,$\sin x$ の原始関数の負号か,部分積分の前の負号を確かめるべきだと分かる.

\begin{example}{部分積分と割り算}
$\int_0^1x\log(1+x)\dd x$ を求めよ.
\end{example}

\paragraphbf{方針}
$\log(1+x)$ を微分し,\,$x$ を積分する.残る有理関数の積分は,分子を分母で割ってから計算する.

\paragraphbf{解答}
$\frac14$

\paragraphbf{解法}
$\log(1+x)$ を微分すると $\frac1{1+x}$ になる.\,$x$ の原始関数は $\frac{x^2}{2}$ だから,残る積分は有理関数になる.そこで
\[
 u=\log(1+x),\qquad v'=x
\]
と選ぶと,
\[
 \int_0^1x\log(1+x)\dd x
 =\left[\frac{x^2}{2}\log(1+x)\right]_0^1
  -\frac12\int_0^1\frac{x^2}{1+x}\dd x
\]
を得る.第1項は $\frac12\log2$ である.第2項では,分子の次数が分母より高いので,\chapterref{章:基本積分}で学んだ割り算を使う.
\[
 x^2=(1+x)(x-1)+1
\]
より,
\[
 \frac{x^2}{1+x}=x-1+\frac1{1+x}
\]
である.したがって
\[
\begin{aligned}
 \int_0^1\frac{x^2}{1+x}\dd x
 &=\left[\frac{x^2}{2}-x+\log(1+x)\right]_0^1\\
 &=-\frac12+\log2
\end{aligned}
\]
となる.これを戻すと
\[
\begin{aligned}
 \int_0^1x\log(1+x)\dd x
 &=\frac12\log2-\frac12\left(-\frac12+\log2\right)\\
 &=\frac14
\end{aligned}
\]
を得る.

\paragraphbf{消えた対数}
部分積分の直後には端点の項に対数が残り,残りの積分にも対数が現れたが,最後に差を取ると打ち消し合った.途中の式の見た目だけで,答えにも対数が必ず残ると決めることはできないのである.
%
\section{加速度と位置}\label{節:加速度と位置}
\chapterref{章:導入}では,速度を短い時間ごとに足して位置の変化を求めた.部分積分を使うと,同じ変化を加速度から表すこともできる.

速度$v(t)$は連続微分可能とし,加速度を$a(t)=v'(t)$とする.\,$a(t)$は速度がどの程度の割合で変わるかを表す量である.時刻$0$から$T$までの位置の変化を$D$とすると,
\[
 D=\int_0^T v(t)\dd t.
\]
ここで$t$の代わりに$t-T$を微分された因子と見て,部分積分してみよう.
\begin{align*}
 D&=[(t-T)v(t)]_0^T-\int_0^T(t-T)v'(t)\dd t\\
 &=Tv(0)+\int_0^T(T-t)a(t)\dd t.
\end{align*}
最初の項は,最初の速度$v(0)$のまま動いた場合の変化である.残りの項は,その後に速度が変わった分を表す.時刻$t$付近で速度が$a(t)\Delta t$だけ変われば,その変化は残りの時間$T-t$にわたって位置へ影響する.そのため,\,$T-t$が掛かると読むこともできる.

例えば加速度が一定値$a$なら,
\[
 D=Tv(0)+a\int_0^T(T-t)\dd t
 =Tv(0)+\frac12aT^2.
\]
導入での$v(t)=2t$は,\,$v(0)=0$,\,$a=2$の場合なので,\,$D=T^2$が再び得られる.最初の長方形の和,速度のグラフの三角形,部分積分という三つの計算が,同じ位置の変化を求めている.

この式は,単に物理の公式を増やすためのものではない.積分する関数を$v$から$v'$へ替えたことで,初めの値と,途中で変わった分とに分けて表せたのである.加速度が与えられている問題では,速度を途中で求めずに使えることもある.

%
\section*{章末問題}
\begin{enumerate}[label=\textbf{\thechapter.\arabic*},leftmargin=*]
 \item\label{問:部分一次指数} $\int x\e^{3x}\dd x$ を求めよ.微分する因子と積分する因子を示し,答えを微分して確かめよ.
 \item\label{問:部分整式対数} $x>0$ において,\,$\int x^2\log x\dd x$ を求めよ.\,$x^2$ と $\log x$ のどちらを微分する側に選ぶか,その理由も述べよ.
 \item\label{問:部分反復} $\int x^3\sin x\dd x$ を求めよ.部分積分を繰り返すたびに,整式の次数がどう変わるかを確かめよ.
 \item\label{問:部分循環} 部分積分を二回使い,\,$\int\e^{2x}\cos x\dd x$ を求めよ.もとの積分が戻った後の等式も書け.
 \item\label{問:部分対数余弦} $x>0$ において,\,$\int\cos(\log x)\dd x$ を求めよ.置換の際には,\,$\dd x$ が何に変わるかも書け.
 \item\label{問:部分置換定積分} $\int_0^1x\log(1+x^2)\dd x$ を求めよ.置換の後に,どの積分を部分積分で求めればよいかを示せ.
 \item\label{問:部分根号定積分} 部分積分を使い,\,$\int_0^1\sqrt{4+x^2}\dd x$ を求めよ.必要になる根号の逆数の原始関数は,微分によって確かめて用いよ.
 \item\label{問:部分情報だけの計算} $f$ は $[0,1]$ で連続微分可能とし,\,$f(0)=1$,\,$f(1)=2$,\,$\int_0^1f(x)\dd x=\frac32$ を満たす.\,$f$ の具体的な式を求めずに,\,$\int_0^1xf'(x)\dd x$ と $\int_0^1(1-x)f'(x)\dd x$ を求めよ.
\end{enumerate}

\smallskip\noindent 解答は\pageref{解答:第7章}ページ.

%
